\section{Method}
\label{sec:method}

We train a neural operator \(\mathcal G_\theta:\mathcal A\to\mathcal U\) with ordinary supervised regression and a stochastic local orbit-consistency objective. Each minibatch contains input--solution pairs \((a,u)\). The supervised term is

\begin{figure}[t]
\centering
\resizebox{0.92\linewidth}{!}{%
\begin{tikzpicture}[
  x=1cm,
  y=1cm,
  line cap=round,
  line join=round,
  >=Latex,
  font=\small
]
\definecolor{groupInk}{RGB}{84,117,158}
\definecolor{groupFill}{RGB}{235,241,249}
\definecolor{planeFill}{RGB}{242,246,252}
\definecolor{inputBlue}{RGB}{28,96,170}
\definecolor{inputFill}{RGB}{235,243,253}
\definecolor{outputPurple}{RGB}{105,55,155}
\definecolor{outputFill}{RGB}{243,236,250}
\definecolor{softGray}{RGB}{105,105,105}

\tikzset{
  point/.style={circle, minimum size=4.8pt, inner sep=0pt},
  groupArrow/.style={-{Latex[length=2.3mm,width=1.6mm]}, dashed, draw=groupInk, line width=0.95pt},
  inputArrow/.style={-{Latex[length=2.4mm,width=1.7mm]}, draw=inputBlue, line width=1.05pt},
  outputArrow/.style={-{Latex[length=2.4mm,width=1.7mm]}, draw=outputPurple, line width=1.05pt},
  tangentArrow/.style={-{Latex[length=1.8mm,width=1.25mm]}, line width=0.82pt},
  connector/.style={-{Latex[length=1.9mm,width=1.35mm]}, draw=softGray!70, line width=0.72pt},
  noteBox/.style={rounded corners=5pt, draw=softGray!48, fill=white, line width=0.65pt, align=center, inner xsep=5pt, inner ysep=3pt}
}

% --------------------------------------------------------------------------
% Lie group and tangent space.
% --------------------------------------------------------------------------
\path[
  fill=groupFill,
  draw=groupInk!65,
  line width=0.9pt,
  fill opacity=0.78
]
  (0.35,4.72)
  .. controls (1.10,5.58) and (2.10,5.78) .. (3.05,5.55)
  .. controls (4.20,5.28) and (5.12,5.90) .. (6.18,6.08)
  .. controls (7.60,6.32) and (8.90,5.74) .. (9.78,4.88)
  .. controls (8.80,4.46) and (7.82,4.32) .. (6.72,4.56)
  .. controls (5.70,4.78) and (4.92,4.38) .. (3.90,4.05)
  .. controls (2.70,3.67) and (1.42,3.88) .. (0.35,4.72)
  -- cycle;
\node[groupInk, font=\large] at (1.32,4.98) {$G$};

\path[
  fill=planeFill,
  draw=inputBlue!62,
  line width=0.75pt,
  fill opacity=0.82
]
  (2.22,6.06) -- (7.28,6.06) -- (8.06,7.92) -- (3.02,7.92) -- cycle;
\node[font=\normalsize, anchor=west] at (3.02,7.64) {$\mathfrak g=T_eG$};

\coordinate (e) at (5.34,6.28);
\coordinate (epsx) at (6.18,6.98);
\coordinate (gpt) at (8.20,4.96);
\node[point, fill=inputBlue] at (e) {};
\node[point, fill=inputBlue!82] at (epsx) {};
\node[point, fill=inputBlue] at (gpt) {};
\node[font=\normalsize, anchor=north] at ($(e)+(0,-0.12)$) {$e$};

\draw[inputBlue, -{Latex[length=2.4mm,width=1.65mm]}, line width=1.05pt]
  (e) -- ($(e)+(0.00,1.32)$) node[above, font=\normalsize] {$X$};
\draw[inputBlue, -{Latex[length=2.4mm,width=1.65mm]}, line width=1.05pt]
  (e) -- ($(e)+(-1.08,0.88)$) node[left=2pt, font=\normalsize] {$X_1$};
\draw[inputBlue, -{Latex[length=2.4mm,width=1.65mm]}, line width=1.05pt]
  (e) -- ($(e)+(1.42,1.02)$) node[above right=2pt, font=\normalsize] {$X_2$};

\node[
  inputBlue,
  font=\scriptsize,
  anchor=north east,
  fill=white,
  fill opacity=0.82,
  text opacity=1,
  inner sep=1pt
] at ($(epsx)+(-0.07,-0.08)$)
  {$\epsilon X\in\mathfrak g$};
\draw[groupArrow]
  (epsx)
  .. controls (7.05,6.98) and (8.02,6.18) ..
  node[above right=0pt, font=\footnotesize, text=groupInk, pos=0.76]
    {$g=\exp(\epsilon X)\in G$}
  (gpt);

% Bridge from the sampled group element to induced field actions.
\node[noteBox, font=\footnotesize] (bridge) at (5.72,3.62)
  {$g\in G$ induces\\[-0.5mm]$(T_g^{\mathcal A},\,T_g^{\mathcal U})$};
\draw[connector] (gpt) .. controls (7.58,4.36) and (6.74,4.02) .. (bridge.north);

% --------------------------------------------------------------------------
% Input and output field-space orbits.
% --------------------------------------------------------------------------
\coordinate (a) at (2.18,2.02);
\coordinate (ag) at (7.78,2.34);
\coordinate (u) at (2.18,0.66);
\coordinate (ug) at (7.78,0.86);

\path[
  fill=inputFill,
  draw=inputBlue!38,
  line width=0.75pt,
  fill opacity=0.55
]
  (0.70,1.72)
  .. controls (1.45,2.42) and (2.46,2.62) .. (3.42,2.42)
  .. controls (4.75,2.14) and (5.82,2.72) .. (7.08,2.78)
  .. controls (8.18,2.84) and (9.02,2.55) .. (9.64,2.06)
  .. controls (8.88,1.72) and (7.98,1.56) .. (6.86,1.72)
  .. controls (5.64,1.90) and (4.74,1.60) .. (3.52,1.40)
  .. controls (2.36,1.20) and (1.42,1.34) .. (0.70,1.72)
  -- cycle;

\path[
  fill=outputFill,
  draw=outputPurple!40,
  line width=0.75pt,
  fill opacity=0.58
]
  (0.70,0.32)
  .. controls (1.56,0.96) and (2.46,1.02) .. (3.44,0.92)
  .. controls (4.70,0.78) and (5.78,1.18) .. (7.04,1.18)
  .. controls (8.10,1.18) and (9.06,0.98) .. (9.62,0.50)
  .. controls (8.88,0.18) and (7.94,0.10) .. (6.82,0.22)
  .. controls (5.64,0.34) and (4.68,0.12) .. (3.52,0.04)
  .. controls (2.32,-0.04) and (1.34,0.06) .. (0.70,0.32)
  -- cycle;

\node[inputBlue, align=left, font=\footnotesize, anchor=west] at (0.70,3.05)
  {$\mathcal O_a^{\mathcal A}=\{T_h^{\mathcal A}a:h\in G\}\subset\mathcal A$};
\node[outputPurple, align=left, font=\footnotesize, anchor=west] at (0.70,-0.34)
  {$\mathcal O_u^{\mathcal U}=\{T_h^{\mathcal U}u:h\in G\}\subset\mathcal U$};

\draw[inputArrow]
  (a) .. controls (3.65,2.44) and (5.80,2.58) ..
  node[above=2pt, font=\footnotesize, text=inputBlue, pos=0.55] {$T_g^{\mathcal A}$}
  (ag);
\draw[outputArrow]
  (u) .. controls (3.60,0.92) and (5.82,1.06) ..
  node[below=2pt, font=\footnotesize, text=outputPurple, pos=0.55] {$T_g^{\mathcal U}$}
  (ug);

\draw[tangentArrow, draw=inputBlue!80]
  ($(a)+(0.28,0.08)$) .. controls (2.74,2.18) and (3.02,2.22) .. (3.34,2.23);
\node[inputBlue, font=\scriptsize, anchor=south] at (2.94,2.28) {$X^{\mathcal A}a$};
\draw[tangentArrow, draw=outputPurple!82]
  ($(u)+(0.28,0.04)$) .. controls (2.72,0.78) and (3.02,0.80) .. (3.34,0.80);
\node[outputPurple, font=\scriptsize, anchor=north] at (2.95,0.55) {$X^{\mathcal U}u$};

\node[point, fill=inputBlue] at (a) {};
\node[point, fill=inputBlue] at (ag) {};
\node[point, fill=outputPurple] at (u) {};
\node[point, fill=outputPurple] at (ug) {};

\node[inputBlue, font=\normalsize, anchor=south] at ($(a)+(0,0.12)$) {$a$};
\node[inputBlue, font=\normalsize, anchor=south] at ($(ag)+(0,0.12)$) {$T_g^{\mathcal A}a$};
\node[outputPurple, font=\normalsize, anchor=north] at ($(u)+(0,-0.12)$) {$u$};
\node[outputPurple, font=\normalsize, anchor=north] at ($(ug)+(0,-0.12)$) {$T_g^{\mathcal U}u$};

\draw[connector] (bridge.south west) .. controls (4.72,3.00) and (3.40,2.74) .. (3.18,2.38);
\draw[connector] (bridge.south east) .. controls (6.62,2.70) and (6.94,1.58) .. (6.86,1.02);

% Exact solution map and the clean equivariance relation.
\draw[-{Latex[length=2.3mm,width=1.6mm]}, draw=black!78, line width=0.95pt]
  (a) -- node[left=3pt, font=\footnotesize] {$\mathcal S$} (u);
\draw[-{Latex[length=2.0mm,width=1.4mm]}, draw=black!45, line width=0.72pt, dashed]
  (ag) -- node[right=3pt, font=\scriptsize, text=black!65] {$\mathcal S$} (ug);
\node[font=\footnotesize, align=center, text=black!72] at (5.45,-0.05)
  {$\mathcal S(T_g^{\mathcal A}a)=T_g^{\mathcal U}\mathcal S(a)$};
\end{tikzpicture}

}
\caption{
Lie-group actions on input and output field spaces.
A local Lie-algebra step \(\epsilon X\in\mathfrak g=T_eG\) maps through the exponential to a group element \(g=\exp(\epsilon X)\in G\). The same \(g\) induces paired actions \(T_g^{\mathcal A}\) and \(T_g^{\mathcal U}\) on the input and output orbits. For an equivariant exact solution operator \(\mathcal S\), the diagram commutes: \(\mathcal S(T_g^{\mathcal A}a)=T_g^{\mathcal U}\mathcal S(a)\).
}
\label{fig:lie-group-actions}
\end{figure}

\begin{equation}
\mathcal L_{\mathrm{sup}}(\theta)=\mathbb E_{(a,u)}\left[
\frac{\|\mathcal G_\theta(a)-u\|_2}{\max(\|u\|_2,10^{-8})}
\right].
\end{equation}
For a sampled Lie-algebra direction and finite transform \(g\), let
\(\epsilon\) denote the step magnitude (for the two-dimensional Galilean action,
\(\epsilon=\max(\|\delta\|_2,10^{-6})\)).  The implemented local
orbit-consistency term is
\begin{equation}
\mathcal L_{\mathrm{orb}}(\theta)=\mathbb E_{a,g}\left[
\frac{\operatorname{MSE}\!\left(
\mathcal G_\theta(T_g^{\mathcal A}a)-T_g^{\mathcal U}\mathcal G_\theta(a)
\right)}{\epsilon^2+\eta}\right],
\end{equation}
where MSE averages over all spatial and output entries of each example before the
batch average.  Supervised Lie augmentation uses the same unsquared per-example
relative loss,
\begin{equation}
\mathcal L_{\mathrm{aug}}(\theta)=\mathbb E_{(a,u),g}\left[
\frac{\|\mathcal G_\theta(T_g^{\mathcal A}a)-T_g^{\mathcal U}u\|_2}
{\max(\|T_g^{\mathcal U}u\|_2,10^{-8})}\right].
\end{equation}
The combined method minimizes
\(\mathcal L_{\mathrm{sup}}+\lambda_{\mathrm{aug}}\mathcal L_{\mathrm{aug}}
+\lambda_{\mathrm{orb}}\mathcal L_{\mathrm{orb}}\); terms not used by a method
are omitted.

\paragraph{What locality and step normalization add.}
A generic finite-transform consistency penalty would instead minimize
\(\mathbb E_{a,g}[\operatorname{MSE}(\Delta_\theta(a,g))]\), without dividing
by the local step.  Near the identity,
\(\Delta_\theta(a,\exp(\epsilon X))=\epsilon d_\theta(a;X)+O(\epsilon^2)\).
The unnormalized penalty therefore scales as
\(\mathbb E[\epsilon^2\lVert d_\theta\rVert^2]+O(\mathbb E|\epsilon|^3)\):
larger sampled steps receive quadratically greater leading-order weight, and the
loss scale changes when the transform radius changes.  Dividing by
\(\epsilon^2+\eta\) removes this leading scale when \(\epsilon^2\gg\eta\), so
the weight is interpretable against a tangent defect and is less mechanically
coupled to the chosen local radius.  The stabilizer caps the weight near zero.
This normalization does not discard finite-orbit information: at nonzero steps
the exact nonlinear group action retains higher-order terms that tangent
propagation at the identity does not see.  Section~\ref{sec:experiments} compares
the normalized objective with both explicit tangent propagation and a
scale-matched unnormalized finite-transform control.

\begin{figure}[t]
\centering

\definecolor{locoBlue}{RGB}{28,96,170}
\definecolor{locoPurple}{RGB}{105,55,155}
\definecolor{locoRed}{RGB}{190,45,40}

\resizebox{0.88\linewidth}{!}{%
\begin{tikzpicture}[
  >=Latex,
  font=\small,
  action/.style={-{Latex[length=2.5mm,width=1.8mm]}, line width=.95pt},
  model/.style={-{Latex[length=2.5mm,width=1.8mm]}, line width=.95pt, draw=black!76},
  defect/.style={{Latex[length=2mm,width=1.5mm]}-{Latex[length=2mm,width=1.5mm]}, dashed, line width=.85pt, draw=locoRed},
  pt/.style={circle, minimum size=4.5pt, inner sep=0pt},
  lab/.style={font=\footnotesize, inner sep=1pt}
]

% Coordinates.
\coordinate (a)      at (0,1.26);
\coordinate (ag)     at (2.12,1.36);
\coordinate (u)      at (0,0);
\coordinate (ubar)   at (1.66,0.56);
\coordinate (utilde) at (2.52,-0.42);

% Faint orbit guides.
\draw[locoBlue!24, line width=1.0pt]
  (-.32,1.16) .. controls (-.16,1.22) and (-.08,1.24) .. (a)
  .. controls (.58,1.58) and (1.40,1.62) .. (ag)
  .. controls (2.34,1.38) and (2.58,1.38) .. (2.78,1.35);
\draw[locoPurple!24, line width=1.0pt]
  (-.32,-.18) .. controls (-.16,-.10) and (-.08,-.04) .. (u)
  .. controls (.52,.42) and (1.08,.62) .. (ubar)
  .. controls (2.02,.46) and (2.42,.12) .. (2.78,.02);

% Two paths being compared.
\draw[action, draw=locoBlue]
  (a) .. controls (.58,1.58) and (1.40,1.62) ..
  node[lab, above=2pt, text=locoBlue, pos=.52] {$T_g^{\mathcal A}$} (ag);

\draw[model]
  (a) -- node[lab, left=3pt] {$G_\theta$} (u);

\draw[action, draw=locoPurple]
  (u) .. controls (.52,.42) and (1.08,.62) ..
  node[lab, below=2pt, text=locoPurple, pos=.52] {$T_g^{\mathcal U}$} (ubar);

\draw[model]
  (ag) -- node[lab, right=3pt] {$G_\theta$} (utilde);

% Points.
\node[pt, fill=locoBlue] at (a) {};
\node[pt, fill=locoBlue] at (ag) {};
\node[pt, fill=locoPurple] at (u) {};
\node[pt, fill=locoPurple] at (ubar) {};
\node[pt, fill=white, draw=locoPurple, line width=.75pt] at (utilde) {};

% Short labels.
\node[lab, text=locoBlue, above left=1pt] at (a) {$a$};
\node[lab, text=locoBlue, above right=1pt] at (ag) {$a_g$};

\node[lab, text=locoPurple, below left=1pt] at (u) {$\hat u$};
\node[lab, text=locoPurple, above left=2pt] at (ubar) {$\bar u_g$};
\node[lab, text=locoPurple, below right=2pt] at (utilde) {$\tilde u_g$};

% Equivariance defect.
\draw[defect] ($(ubar)+(0.06,-0.025)$) -- ($(utilde)+(-0.06,0.025)$);
\node[lab, text=locoRed, anchor=north east] at ($(ubar)!0.5!(utilde)+(-0.035,-0.055)$)
  {$\Delta_\theta(a,g)$};

\end{tikzpicture}%
}

\caption{
Local orbit consistency.
For \(g=\exp(\epsilon X)\), let
\(a_g=T_g^{\mathcal A}a\),
\(\hat u=G_\theta(a)\),
\(\bar u_g=T_g^{\mathcal U}\hat u\), and
\(\tilde u_g=G_\theta(a_g)\).
LOCO penalizes the red equivariance defect
\(\Delta_\theta(a,g)=\tilde u_g-\bar u_g\), equivalently
\(\operatorname{MSE}(\tilde u_g-\bar u_g)/(\epsilon^2+\eta)\) in the
discrete implementation.
}
\label{fig:loco-intuition}
\end{figure}

\begin{figure*}[t]
\centering
\begin{tikzpicture}[
  x=1cm,
  y=1cm,
  line cap=round,
  line join=round,
  >=Latex,
  font=\small
]
\definecolor{orbitBlue}{RGB}{12,67,143}
\definecolor{orbitBlueLight}{RGB}{219,231,247}
\definecolor{orbitPurple}{RGB}{87,31,143}
\definecolor{orbitPurpleLight}{RGB}{235,223,246}
\definecolor{defectRed}{RGB}{190,27,18}
\definecolor{softGray}{RGB}{95,95,95}

% Titles.
\node[orbitBlue, font=\bfseries\large] at (3.25,9.60) {Input orbit};
\node[orbitBlue, font=\large] at (3.25,9.06)
  {$\mathcal O_{\mathcal A}=\{T_g^{\mathcal A}a\mid g\in\mathcal G\}$};
\node[orbitPurple, font=\bfseries\large] at (12.55,9.60) {Output orbit};
\node[orbitPurple, font=\large] at (12.55,9.06)
  {$\mathcal O_{\mathcal U}=\{T_g^{\mathcal U}u\mid g\in\mathcal G\}$};

% Input orbit sheet.
\path[
  fill=orbitBlueLight,
  draw=orbitBlue!45,
  line width=1pt,
  fill opacity=0.62
]
  (0.20,5.35)
  .. controls (0.55,6.15) and (0.95,6.70) .. (1.70,6.80)
  .. controls (2.95,6.95) and (3.60,7.45) .. (4.30,8.15)
  .. controls (5.25,9.10) and (6.40,8.92) .. (6.95,8.18)
  .. controls (7.35,7.65) and (8.15,7.95) .. (8.10,7.18)
  .. controls (8.03,6.40) and (7.25,6.35) .. (7.05,5.65)
  .. controls (6.82,4.85) and (6.05,4.55) .. (5.05,4.70)
  .. controls (4.00,4.86) and (3.20,4.42) .. (2.35,3.95)
  .. controls (1.50,3.48) and (0.78,3.65) .. (0.32,4.28)
  .. controls (0.10,4.58) and (0.02,4.95) .. (0.20,5.35)
  -- cycle;

\draw[orbitBlue!55, dashed, line width=1.1pt]
  (0.62,4.18)
  .. controls (1.05,4.98) and (1.55,5.65) .. (2.08,6.00)
  .. controls (2.96,6.58) and (3.88,6.94) .. (4.84,7.10)
  .. controls (5.55,7.22) and (6.18,7.74) .. (6.78,7.74);
\draw[orbitBlue, line width=1.25pt]
  (2.08,6.00) .. controls (2.75,6.30) and (3.60,6.68) .. (4.84,7.10);
\draw[orbitBlue, -{Latex[length=3mm]}, line width=1.65pt]
  (2.08,6.00) .. controls (2.82,6.50) and (3.76,6.90) .. (4.84,7.10);
\draw[orbitBlue!85, -{Latex[length=2.4mm]}, line width=0.85pt]
  (2.08,6.00) .. controls (3.15,5.92) and (4.16,6.44) .. (4.84,7.10);

\path[fill=orbitBlue, draw=black!45, line width=0.5pt] (2.08,6.00) circle (0.12);
\path[fill=white, opacity=0.75] (2.04,6.05) circle (0.035);
\path[fill=orbitBlue, draw=black!45, line width=0.5pt] (4.84,7.10) circle (0.12);
\path[fill=white, opacity=0.75] (4.80,7.15) circle (0.035);

\node[orbitBlue, font=\large] at (1.88,5.38) {$a$};
\node[orbitBlue, font=\large] at (4.72,7.52) {$T_g^{\mathcal A}a$};
\node[orbitBlue, align=center, font=\footnotesize, inner sep=1pt]
  at (3.72,5.38)
  {$g=\exp(\epsilon X)$\\[-1mm]{\footnotesize local Lie-group step}};

% Operator arrow.
\draw[black!88, -{Latex[length=4mm,width=3mm]}, line width=1.55pt]
  (7.18,6.46) .. controls (7.92,6.88) and (8.78,6.78) .. (9.72,6.42);
\node[font=\huge, fill=white, fill opacity=0.75, text opacity=1, inner sep=0.5pt]
  at (8.10,7.12) {$G_\theta$};

% Output orbit sheet.
\path[
  fill=orbitPurpleLight,
  draw=orbitPurple!45,
  line width=1pt,
  fill opacity=0.68
]
  (8.15,4.50)
  .. controls (8.50,5.72) and (9.62,6.05) .. (10.15,6.82)
  .. controls (10.98,8.02) and (12.00,8.72) .. (13.40,8.68)
  .. controls (14.70,8.64) and (15.55,8.05) .. (15.82,7.05)
  .. controls (15.22,6.90) and (14.95,6.22) .. (14.88,5.62)
  .. controls (14.75,4.50) and (13.62,4.45) .. (12.64,4.78)
  .. controls (11.72,5.08) and (10.70,4.92) .. (9.75,4.62)
  .. controls (9.15,4.42) and (8.58,4.30) .. (8.15,4.50)
  -- cycle;

\draw[orbitPurple!35, dashed, line width=1.1pt]
  (9.70,5.62)
  .. controls (10.40,6.40) and (11.36,7.12) .. (12.65,7.33)
  .. controls (13.46,7.46) and (14.12,7.58) .. (14.82,7.48);
\draw[orbitPurple, line width=1.2pt]
  (9.70,5.62) .. controls (10.82,6.20) and (11.85,6.85) .. (12.65,7.33);
\draw[orbitPurple, -{Latex[length=3mm]}, line width=1.65pt]
  (9.70,5.62) .. controls (10.58,6.38) and (11.54,6.92) .. (12.65,7.33);
\draw[orbitPurple!75, -{Latex[length=2.3mm]}, line width=0.85pt]
  (9.70,5.62) .. controls (10.80,5.75) and (11.95,5.70) .. (12.95,5.82);
\draw[black!85, densely dashed, line width=1.05pt]
  (12.65,7.33) -- (12.95,5.82);

\path[fill=orbitPurple, draw=black!45, line width=0.5pt] (9.70,5.62) circle (0.12);
\path[fill=white, opacity=0.75] (9.66,5.67) circle (0.035);
\path[fill=orbitPurple, draw=black!45, line width=0.5pt] (12.65,7.33) circle (0.12);
\path[fill=white, opacity=0.75] (12.61,7.38) circle (0.035);
\path[fill=orbitPurple!35, draw=orbitPurple!70, line width=0.6pt] (12.95,5.82) circle (0.12);

\node[orbitPurple, font=\large] at (9.48,5.18) {$G_\theta(a)$};
\node[orbitPurple, font=\large] at (12.38,7.77) {$T_g^{\mathcal U}G_\theta(a)$};
\node[orbitPurple, font=\large] at (13.58,5.30) {$G_\theta(T_g^{\mathcal A}a)$};
\draw[defectRed, -{Latex[length=2.3mm]}, line width=0.9pt]
  (13.18,6.60) -- (12.83,6.55);
\node[defectRed, font=\large] at (13.88,6.58) {$\Delta_\theta(a,g)$};

% Central loss box.
\draw[rounded corners=7pt, draw=black!62, line width=0.9pt, fill=white]
  (4.32,3.16) rectangle (11.68,4.50);
\node[defectRed, font=\large] at (8.00,4.25) {penalize local orbit mismatch};
\node[font=\Large] at (8.00,3.77)
  {$\Delta_\theta(a,g)=\textcolor{orbitPurple}{G_\theta(T_g^{\mathcal A}a)}
  -\textcolor{orbitPurple}{T_g^{\mathcal U}G_\theta(a)}$};
\node[defectRed, font=\footnotesize, fill=white, inner xsep=3pt] at (8.00,3.16)
  {equivariance defect};

% Intuition panel.
\draw[rounded corners=8pt, draw=black!62, line width=0.85pt, fill=white]
  (0.35,0.50) rectangle (15.65,3.03);
\node[font=\normalsize] at (8.00,2.78)
  {Intuition: \textcolor{orbitBlue}{local consistency} $\to$
   \textcolor{orbitPurple}{improved symmetry robustness} on larger orbit shifts};

% Left mini diagram.
\node[orbitBlue, align=center, font=\footnotesize] at (1.62,2.05)
  {small (local) shift\\[-1mm]$g=\exp(\epsilon X)$};
\path[fill=orbitBlueLight, draw=orbitBlue!45, line width=0.6pt, fill opacity=0.65]
  (2.50,1.28)
  .. controls (2.82,1.70) and (3.18,1.68) .. (3.44,1.88)
  .. controls (3.80,2.15) and (4.30,2.14) .. (4.62,1.88)
  .. controls (4.43,1.55) and (4.12,1.38) .. (3.62,1.35)
  .. controls (3.22,1.34) and (2.86,1.12) .. (2.50,1.28) -- cycle;
\draw[orbitBlue!70, dashed, line width=0.7pt]
  (2.72,1.38)
  .. controls (2.88,1.54) and (3.00,1.64) .. (3.15,1.70)
  .. controls (3.46,1.84) and (3.72,1.88) .. (4.03,1.90)
  .. controls (4.18,1.91) and (4.30,1.92) .. (4.42,1.92);
\draw[orbitBlue, -{Latex[length=2mm]}, line width=1pt]
  (3.15,1.70) -- (4.03,1.90);
\path[fill=orbitBlue, draw=black!40, line width=0.35pt] (3.15,1.70) circle (0.075);
\path[fill=orbitBlue, draw=black!40, line width=0.35pt] (4.03,1.90) circle (0.075);

\draw[black, -{Latex[length=2.4mm]}, line width=1pt]
  (4.62,1.88) .. controls (4.92,2.02) and (5.22,2.00) .. (5.58,1.78);
\node at (5.04,2.13) {$G_\theta$};

\path[fill=orbitPurpleLight, draw=orbitPurple!45, line width=0.6pt, fill opacity=0.68]
  (5.48,1.42)
  .. controls (5.71,1.78) and (6.12,2.10) .. (6.66,2.10)
  .. controls (6.94,2.10) and (7.14,1.94) .. (7.24,1.68)
  .. controls (6.88,1.54) and (6.58,1.28) .. (6.12,1.34)
  .. controls (5.78,1.38) and (5.58,1.30) .. (5.48,1.42) -- cycle;
\draw[orbitPurple, -{Latex[length=1.9mm]}, line width=0.95pt]
  (6.01,1.76) -- (6.54,1.92);
\draw[orbitPurple!70, -{Latex[length=1.7mm]}, line width=0.65pt]
  (6.01,1.76) -- (6.66,1.65);
\draw[black!75, densely dashed, line width=0.65pt] (6.54,1.92) -- (6.68,1.65);
\path[fill=orbitPurple, draw=black!40, line width=0.35pt] (6.01,1.76) circle (0.075);
\path[fill=orbitPurple, draw=black!40, line width=0.35pt] (6.54,1.92) circle (0.075);
\path[fill=orbitPurple!35, draw=orbitPurple!70, line width=0.35pt] (6.68,1.65) circle (0.075);
\node[defectRed, font=\large] at (6.38,1.13) {$\approx 0$};

\draw[densely dashed, softGray, line width=0.8pt] (7.42,0.90) -- (7.42,2.38);

% Right mini diagram.
\node[orbitBlue, align=center, font=\footnotesize] at (8.42,2.05)
  {larger shift\\[-1mm]$g=\exp(tX)$\\[-1mm]{\footnotesize(larger $t$)}};
\path[fill=orbitBlueLight, draw=orbitBlue!45, line width=0.6pt, fill opacity=0.65]
  (9.25,1.18)
  .. controls (9.72,1.58) and (10.08,1.68) .. (10.50,1.92)
  .. controls (10.92,2.16) and (11.46,2.16) .. (11.72,1.84)
  .. controls (11.36,1.52) and (10.90,1.34) .. (10.36,1.32)
  .. controls (9.90,1.30) and (9.54,1.02) .. (9.25,1.18) -- cycle;
\draw[orbitBlue!70, dashed, line width=0.7pt]
  (9.45,1.28)
  .. controls (9.58,1.42) and (9.76,1.52) .. (9.95,1.58)
  .. controls (10.35,1.72) and (10.70,1.90) .. (11.02,1.98)
  .. controls (11.22,2.03) and (11.38,2.03) .. (11.52,2.02);
\draw[orbitBlue, -{Latex[length=2mm]}, line width=1pt]
  (9.95,1.58) -- (11.02,1.98);
\path[fill=orbitBlue, draw=black!40, line width=0.35pt] (9.95,1.58) circle (0.075);
\path[fill=orbitBlue, draw=black!40, line width=0.35pt] (11.02,1.98) circle (0.075);

\draw[black, -{Latex[length=2.4mm]}, line width=1pt]
  (11.72,1.84) .. controls (12.08,2.06) and (12.52,2.06) .. (12.96,1.84);
\node at (12.12,2.17) {$G_\theta$};

\path[fill=orbitPurpleLight, draw=orbitPurple!45, line width=0.6pt, fill opacity=0.68]
  (12.88,1.22)
  .. controls (13.12,1.88) and (13.62,2.40) .. (14.18,2.36)
  .. controls (14.76,2.30) and (15.10,1.98) .. (15.32,1.68)
  .. controls (14.76,1.62) and (14.52,1.20) .. (13.98,1.16)
  .. controls (13.54,1.12) and (13.15,1.02) .. (12.88,1.22) -- cycle;
\draw[orbitPurple, -{Latex[length=2mm]}, line width=1.05pt]
  (13.62,1.66) .. controls (14.02,1.92) and (14.35,2.06) .. (14.72,2.14);
\draw[orbitPurple!70, -{Latex[length=1.8mm]}, line width=0.68pt]
  (13.62,1.66) -- (14.78,1.75);
\draw[black!75, densely dashed, line width=0.65pt] (14.72,2.14) -- (14.78,1.75);
\path[fill=orbitPurple, draw=black!40, line width=0.35pt] (13.62,1.66) circle (0.075);
\path[fill=orbitPurple, draw=black!40, line width=0.35pt] (14.72,2.14) circle (0.075);
\path[fill=orbitPurple!35, draw=orbitPurple!70, line width=0.35pt] (14.78,1.75) circle (0.075);
\node[defectRed, align=center] at (14.48,1.14) {small\\[-1mm]mismatch};

\draw[black, -{Latex[length=3mm]}, line width=1.1pt] (2.42,0.90) -- (12.70,0.90);
\node[orbitBlue] at (1.98,0.90) {small};
\node[font=\footnotesize, fill=white, inner ysep=0.2pt] at (7.68,0.68) {orbit shift size};
\node[orbitBlue] at (13.12,0.90) {larger};

% Legend.
\draw[dashed, rounded corners=4pt, draw=black!50, fill=white] (0.88,-0.14) rectangle (15.12,0.31);
\path[fill=orbitBlue, draw=black!45, line width=0.45pt] (1.20,0.085) circle (0.10);
\node[anchor=west, font=\scriptsize] at (1.38,0.085) {input orbit point};
\path[fill=orbitPurple, draw=black!45, line width=0.45pt] (4.12,0.085) circle (0.10);
\node[anchor=west, font=\scriptsize] at (4.30,0.085) {output orbit point};
\path[fill=orbitPurple!35, draw=orbitPurple!70, line width=0.45pt] (7.20,0.085) circle (0.10);
\node[anchor=west, font=\scriptsize] at (7.38,0.085) {mapped point $G_\theta(T_g^{\mathcal A}a)$};
\draw[black!85, densely dashed, line width=1pt] (11.40,0.085) -- (11.82,0.085);
\node[anchor=west, defectRed, font=\scriptsize] at (12.02,0.085) {defect $\Delta_\theta(a,g)$};
\end{tikzpicture}

\caption{Vector schematic of local orbit consistency. A local Lie-group step moves an input \(a\) to \(T_g^{\mathcal A}a\). The training penalty compares the prediction on the transformed input, \(\mathcal G_\theta(T_g^{\mathcal A}a)\), with the transformed original prediction, \(T_g^{\mathcal U}\mathcal G_\theta(a)\), reducing the equivariance defect along the orbit while leaving inference unchanged.}
\label{fig:orbit-consistency-diagram}
\end{figure*}

The loss is architecture-agnostic in the implementation sense: it requires only model evaluations and known input/output actions, not Fourier layers. We test it primarily with FNO backbones, reproduce the main Galilean comparison with a near-parameter-matched DeepONet, and separately use a non-equivariant molecular MLP as a mechanism check. This does not imply backbone-independent effect size. Transformations are applied to physical-grid input fields; the prediction on the transformed input is compared with the transformed original prediction; gradients pass through both model evaluations. Integer periodic translations are tensor rolls, non-integer periodic translations are spectral shifts, and finite square symmetries are exact tensor rotations and flips. For 1D Burgers, the Galilean action on the solution map \(u_0\mapsto u_T\) is \(T_c^{\mathcal A}u_0(x)=u_0(x)+c\) and \(T_c^{\mathcal U}u_T(x)=u_T(x-cT)+c\).

The same loss supports unlabeled adaptation. Given a pretrained operator \(\mathcal G_{\theta_0}\) and target inputs \(a\sim\mathcal D_{\mathrm{target}}\), we minimize \(\mathcal L_{\mathrm{orb}}^{\mathrm{target}}(\theta)+\beta\|\theta-\theta_0\|_2^2\) over a chosen adaptation set \(\Theta_{\mathrm{adapt}}\), such as all parameters, the last operator block, or the projection head. This adaptation objective requires no PDE solves and no target labels.
